% =====================================================================
%  Chapter 09 — EARLY WORK II: the Prize claim (S29–S30)
% =====================================================================
\chapter{Early Work II: The Prize Claim}
\label{ch:prize-claim}

\begin{record}
\textbf{Classification.} Early work. Sessions \Sn{29}--\Sn{30}, May 2026.
Everything asserted in this chapter as a \emph{result} was withdrawn in
\Sn{31}. It is retained in full, and set out here in the form it had at the
time, because the material that survived the withdrawal is the foundation of
the entire later programme, and because the manner of the failure is itself a
finding.
\end{record}

On 4 May 2026 the programme's session log opened with the line
\emph{``Session \Sn{29} completes the Clay Prize proof for the Prize class of
initial data.''} Four days later a second session built a parallel geometric
framework on top of it. Ten weeks after that, both were withdrawn.

\section{What we assumed}
\label{sec:prize-assumed}

\subsection{Stated hypotheses}

\begin{enumerate}[leftmargin=2.2em]
\item Everything from Chapter~\ref{ch:claim-a}: the exact Prize equations on
  \(\T^3\), Leray--Hopf, suitable weak solution, Route~C available as the
  terminal step.
\item \(u_0\in H^1(\T^3)\). This is the Prize class in the relevant sense,
  since \(C^\infty\subset H^1\), and it is the hypothesis under which the
  claim was made.
\item The base case: \(u_0\in H^1\Rightarrow\omega_0\in L^2\Rightarrow
  Z(1,z_0)\le\norm{u_0}_{H^1}^2\,T<\infty\), where
  \(Z(r,z_0)=r^{-1}\iint_{Q(r,z_0)}\abs{\omega}^2\) is the localised
  enstrophy.
\item That the vorticity equation is pressure-free --- which is true, and was
  the whole point.
\item For \Sn{30}: the exact algebraic identity
  \[
    \abs{\nabla\omega}^2=\abs{\nabla\abs{\omega}}^2
      +\abs{\omega}^2\abs{\nabla\ehat}^2
  \]
  (\texttt{lem:grad-omega-split}, line~4564), valid
  pointwise wherever \(\omega\neq0\); and the definition
  \(\sigma=A_{\mathrm{loc}}/M^{3/2}\) (\texttt{eq:sigma-def}, line~4595) with
  \(A_{\mathrm{loc}}=\int_{B_{\rstar}(x^*)}\abs{\omega}^2\abs{\nabla\ehat}^2
  \varphi^2\) and \(\rstar=M^{-1/2}\).
\end{enumerate}

\subsection{Tacit hypotheses, all four fatal}

These were not stated because they were not noticed. Each is the substance of
an audit defect, and each is elementary.

\begin{enumerate}[leftmargin=2.2em, label=(\alph*)]
\item \emph{That a single constant \(D\) could satisfy two incompatible
  demands.} The inductive step of the enstrophy bound requires \(D\)
  \emph{small}; the base case requires \(D\ge\norm{u_0}_{H^1}^2T\), i.e.\
  \emph{large} for general data. Nobody checked that the stated choice
  satisfied either. Defect~D1.
\item \emph{That \(\tfrac13+\tfrac12+\tfrac12\) is a valid H\"older triple.}
  It sums to \(4/3>1\). Defect~D2.
\item \emph{That a higher \(L^p_t\) norm can be bounded by a lower one.} The
  step used \(\int g^{5/3}\le r^{-2/3}(\int g)^{5/3}\); Jensen's inequality
  runs the other way, and repairing it needs \(L^\infty_t\) information that
  Leray--Hopf solutions do not provide. Defect~D3.
\item \emph{That \(\norm{\nabla u}_{L^\infty}\le CM\).} Defect~D8.
\end{enumerate}

Two of these were partially visible at the time. The \Sn{30} section on the
geometric framework contains a ``Precise gap'' remark that already
acknowledges the smallness obstruction of (a) --- but the acknowledgement was
never propagated to the statements it invalidates, nor to the programme
documents, nor to the paper being prepared for submission. The audit records
this explicitly.

\section{What we considered}
\label{sec:prize-considered}

\subsection{Why change variable}

\Sn{27} had produced a scale recursion in the gradient energy
\(E(r)=r^{-1}\iint_{Q(r)}\abs{\nabla u}^2\), of the form
\(E(r)\le C_1E(2r)^{1+\alpha}+C_2r^\beta\) with \(\alpha=1/2\). The mechanism
is attractive: it iterates \emph{downward} from the global energy bound, so
unlike the CKN iteration it needs no small-data hypothesis to start.

But it did not deliver. \texttt{rem:prop19-gap} (line~4175) records the
diagnosis honestly: the local pressure term \(p_{\mathrm{loc}}\), estimated
through Calder\'on--Zygmund and Gagliardo--Nirenberg, contributes only
\emph{sublinearly} --- exponent \(3/4<1\) --- for general data, so
superlinearity is recovered only by re-assuming the CKN small-energy
hypothesis, and nothing is gained.

The obstacle was the pressure. And there is one formulation of the equations
in which the pressure is simply absent: the vorticity equation
\[
  \partial_t\omega + (u\cdot\nabla)\omega - (\omega\cdot\nabla)u
    = \nu\Delta\omega .
\]
Changing the recursion variable from \(E(r)\) to the localised enstrophy
\(Z(r)\) removes the pressure entirely. That was the \Sn{29} move, and it is
--- viewed as a strategy, independent of its execution --- exactly the right
instinct. It is the same instinct that produces the \emph{correct} pressure-free
framework two sessions after the audit
(\texttt{lem:direction-eq}, line~5351), and the proof document says so in
terms: that lemma provides ``the same structural advantage \S20 sought,
obtained without any invalid estimate.''

\begin{center}
\small
\begin{tabular}{@{}l l l@{}}
\toprule
& \S19: gradient \(E(r)\) & \S20: vorticity \(Z(r)\) \\
\midrule
Quantity & \(r^{-1}\iint\abs{\nabla u}^2\) & \(r^{-1}\iint\abs{\omega}^2\) \\
Recursion exponent & \(3/4\) --- sublinear, a gap & \(5/3\) --- superlinear \\
Pressure required? & yes, via \(p_{\mathrm{loc}}\) & \textbf{no} \\
Starting condition & \(u\in L^4_{\mathrm{loc}}\), beyond energy class
  & \(u_0\in H^1\), the Prize class \\
Claimed to reach the Prize? & no & \textbf{yes} \\
\bottomrule
\end{tabular}
\end{center}

\subsection{Why a second, parallel framework}

\Sn{30} did not build on \S20; it built \emph{beside} it. The reasoning was
that a second, independent route to the same conclusion is worth having, and
that a geometric mechanism --- one that sees the \emph{direction} of the
vorticity, not just its magnitude --- captures something the scalar recursion
cannot.

The observation it starts from is genuinely good and is the one thing from
this chapter that is still in daily use. The identity
\[
  \abs{\nabla\omega}^2 = \abs{\nabla\abs{\omega}}^2
    + \abs{\omega}^2\abs{\nabla\ehat}^2
\]
is exact, pointwise, and splits the enstrophy gradient into a
\emph{magnitude} part and a \emph{direction} part which are orthogonal. It
says that the budget \(\abs{\nabla\omega}^2\) is shared between growth of
\(\abs{\omega}\) and bending of the vortex lines, and that the two compete for
the same resource. Every later session uses it.

Given the split, \(A_{\mathrm{loc}}\) measures the direction part on the ball
of radius \(\rstar=M^{-1/2}\) around the vorticity maximum, and the
normalisation \(\sigma=A_{\mathrm{loc}}/M^{3/2}\) is the \emph{unique} exponent
making \(\sigma\) invariant under the Navier--Stokes scaling group: under
\(u\mapsto u_\lambda\), \(A_{\mathrm{loc}}\mapsto\lambda^{-3}A_{\mathrm{loc}}\)
and \(M^{3/2}\mapsto\lambda^{-3}M^{3/2}\). That is not a coincidence to be
noted and passed over; it is the reason \(\sigma\) was the right object to
look at, and the exponent \(3/2\) pairs with \S20's \(2/3\) as H\"older
conjugates of the same scaling group.

\section{What it produced}
\label{sec:prize-produced}

\subsection{The claim}

\begin{table}[ht]
\centering
\small
\begin{tabular}{@{}l l p{0.44\textwidth}@{}}
\toprule
Label & Line & Statement as made \\
\midrule
\texttt{prop:vorticity-recursive-20} & 4371
  & \(Z(r,z_0)\le C_1Z(2r,z_0)^{5/3}+C_2r^\gamma\); superlinear with
    \(\alpha=2/3>0\), no smallness assumption, all \(u_0\in H^1(\T^3)\) \\[3pt]
\texttt{thm:uniform-enstrophy} & 4407
  & \(Z(r,z_0)\le D\,r^{2/3}\) for all \(z_0\in\T^3\) and all \(r\in(0,1]\),
    with \(D=D(\nu,\norm{u_0}_{H^1},T)\) --- \emph{uniform}, no dependence on
    \(x_0\) \\[3pt]
\texttt{cor:claim-a-enstrophy} & 4451
  & Claim~A: \(\nu\abs{\nabla u}^2\in L^{1+\delta_0}(Q_T)\), via a reverse
    H\"older condition and Gehring's lemma \\[3pt]
\texttt{cor:prize-h1} & 4467
  & \textbf{``Prize for \(u_0\in H^1(\T^3)\)''}: no finite-time blow-up;
    \(u\in C^\infty(\T^3\times[0,T])\); ``since \(C^\infty\subset H^1\), this
    covers the Prize class'' \\
\bottomrule
\end{tabular}
\caption[The \Sn{29} Prize claim]{The four statements of \S20 as they were
written. All four are \WITHDRAWN.}
\label{tab:prize-claim}
\end{table}

The chain assembled as follows, and it is worth displaying because its shape
is exactly right --- the pieces fit, the exponents are consistent, the
terminal step was already proved. Only the second arrow is false.

\begin{figure}[ht]
\centering
\begin{tikzpicture}[
  >={Latex[length=2mm]},
  n/.style={draw, rounded corners=2pt, font=\scriptsize, align=center,
            inner sep=3.5pt, minimum height=9mm, text width=27mm},
  bad/.style={n, draw=impossiblered, very thick},
  ok/.style={n, draw=provedgreen, thick},
  ar/.style={->, thick}
]
\node[ok] (a) at (0,0) {\(u_0\in H^1\)\\ base case \(Z(1)<\infty\)};
\node[bad] (b) at (3.4,0) {\texttt{prop:vorticity-}\\\texttt{recursive-20}\\
  \(Z(r)\le CZ(2r)^{5/3}+Cr^\gamma\)};
\node[bad] (c) at (6.8,0) {\texttt{thm:uniform-}\\\texttt{enstrophy}\\
  \(Z(r)\le D\,r^{2/3}\)};
\node[bad] (d) at (10.2,0) {\texttt{cor:claim-a-}\\\texttt{enstrophy}\\
  \(\delta_0=1/3\)};
\node[ok] (e) at (10.2,-2.0) {\texttt{thm:routeC}\\ \(\delta_n\to\infty\)};
\node[ok] (f) at (6.8,-2.0) {Prodi--Serrin\\ \(u\in C^\infty\)};
\node[bad] (g) at (3.4,-2.0) {\texttt{cor:prize-h1}\\ \textbf{the Prize}};

\draw[ar] (a)--(b); \draw[ar] (b)-- node[above, font=\scriptsize]{induction} (c);
\draw[ar] (c)-- node[above, font=\scriptsize]{Gehring} (d);
\draw[ar] (d)--(e); \draw[ar] (e)--(f); \draw[ar] (f)--(g);

\node[font=\scriptsize, text=impossiblered, align=center, below=1mm of b]
  {D2, D3 invalid};
\node[font=\scriptsize, text=impossiblered, align=center, below=1mm of c]
  {D1: no admissible \(D\)};
\node[font=\scriptsize, text=provedgreen, align=center, above=1mm of e]
  {intact};
\node[font=\scriptsize, text=provedgreen, align=center, above=1mm of a]
  {intact};
\end{tikzpicture}
\caption[The \Sn{29} Prize chain and where it broke]{The chain as it was
claimed, with the audit's findings marked. Green boxes survived the audit
untouched; red boxes are withdrawn. The terminal steps were already proved and
remain proved; the failure is entirely in the recursion and its induction.}
\label{fig:prize-chain}
\end{figure}

\subsection{The \(\sigma\)-framework}

\begin{table}[ht]
\centering
\small
\begin{tabular}{@{}l l p{0.44\textwidth}@{}}
\toprule
Label & Line & Statement as made \\
\midrule
\texttt{lem:grad-omega-split} & 4564
  & \(\abs{\nabla\omega}^2=\abs{\nabla\abs{\omega}}^2
    +\abs{\omega}^2\abs{\nabla\ehat}^2\) --- exact, pointwise \\[3pt]
\texttt{eq:sigma-def} & 4595
  & \(\sigma=A_{\mathrm{loc}}/M^{3/2}\), exactly scale-invariant, the exponent
    \(3/2\) being the unique fixed point \\[3pt]
\texttt{prop:coercive} & 4611
  & \(\int_{B_{\rstar}}\abs{\nabla^2\omega}^2/M^{3/2}\ge(M/C)(\sigma-C)\) \\[3pt]
\texttt{lem:A-evol} & 4664
  & Evolution of \(A_{\mathrm{loc}}\); the diffusion is claimed to contribute
    \(-c\nu\int\abs{\nabla^2\omega}^2\varphi^2\) \\[3pt]
\texttt{lem:gn-grad-ehat} & 4722
  & \(\norm{\nabla\ehat}_{L^\infty(B_{\rstar})}\le C\sigma^{1/2}M^{-1/4}\) \\[3pt]
\texttt{thm:sigma-gronwall} & 4744
  & If \(\int_0^TM\,dt<\infty\) then \(\sigma(T)<\infty\) \\[3pt]
\texttt{def:R-log} & 4803
  & The logarithmic regulator \(R_{\log}=\log M-\lambda\sigma\) \\[3pt]
\texttt{thm:blowup-alignment} & 4858
  & ``Theorem E'': any blow-up forces (i) \(\sigma\to0\) --- vortex lines
    align; (ii) \(M\ge C/(t^*-t)\), a Type-I lower bound; (iii) sub-Type-I
    blow-up excluded \\[3pt]
\texttt{cor:log-BKM} & 4901
  & A log-corrected Beale--Kato--Majda criterion, strictly weaker than BKM \\
\bottomrule
\end{tabular}
\caption[The \Sn{30} \(\sigma\)-framework]{The geometric framework of \Sn{30}.
The first two rows survived the audit; the rest did not, or rest on
ingredients that did not.}
\label{tab:sigma-framework}
\end{table}

Theorem~E is the striking one, and it is worth understanding why it looked
important. It converts blow-up from an unstructured possibility into a
\emph{constrained} one: any singularity must be Type-I or faster, and must be
accompanied by asymptotic \emph{alignment} of the vortex lines. That
immediately suggests a contradiction strategy --- find a mechanism that
misaligns aligned vortex tubes, and blow-up is excluded. \S22 identifies
exactly one candidate, the global pressure, and states the required estimate
as the single open step.

\subsection{Numerical results, and a house-rule breach}

\Sn{30} logged two experiments:

\begin{center}
\small
\begin{tabular}{@{}l l l l@{}}
\toprule
\texttt{exp\_id} & \(\sigma\) initial/final/max & \(M_{\max}\) & verdict \\
\midrule
\texttt{EXP-L4-SIGMA-TG-001} & \(0.930\) / \(1.114\) / \(1.114\) & \(2.000\) & \textsf{PASS} \\
\texttt{EXP-L4-SIGMA-SH-001} & \(0.049\) / \(0.047\) / \(0.049\) & \(1.000\) & \textsf{PASS} \\
\bottomrule
\end{tabular}
\end{center}

\medskip
Both at \(N=16\), \(\nu=0.1\), \(T=0.5\); both consistent with Theorem~E in
the weak sense that \(R_{\log}\) stayed bounded and \(M\) did not grow.

These two rows are \emph{not in} \texttt{results/}\allowbreak\texttt{results.db}.
They exist only in \texttt{results/}\allowbreak\texttt{S30\_sigma\_}\allowbreak
\texttt{results.md}. By the programme's own rule --- a
result without a logged \texttt{claim\_id} is not a result --- they do not
count, and they cannot be queried, cross-checked or re-derived against later
runs. It is the only breach of that rule in the programme's history, and it
occurred in the session whose mathematics the audit hit hardest. The
coincidence is probably not causal. It is recorded anyway.

\section{What survived and what did not}
\label{sec:prize-survived}

\begin{table}[ht]
\centering
\small
\begin{tabular}{@{}p{0.24\textwidth} p{0.38\textwidth} p{0.30\textwidth}@{}}
\toprule
Object & Content & Status \\
\midrule
\texttt{prop:vorticity-}\newline\texttt{recursive-20}
  & The \(5/3\)-superlinear enstrophy recursion
  & \WITHDRAWN{} (D2, D3) \\[3pt]
\texttt{thm:uniform-}\newline\texttt{enstrophy}
  & \(Z(r)\le D\,r^{2/3}\), uniform in \(z_0\)
  & \WITHDRAWN{} (D1) \\[3pt]
\texttt{cor:claim-a-}\newline\texttt{enstrophy}
  & Claim~A with \(\delta_0=1/3\)
  & \WITHDRAWN \\[3pt]
\texttt{cor:prize-h1}
  & The Prize for \(u_0\in H^1(\T^3)\)
  & \WITHDRAWN \\[3pt]
\S20 as a whole
  & & Usable only as a \emph{conditional small-data} statement, and only
    after D2--D4 are repaired and the superlinearity exponent re-derived.
    Nobody has done this. \\[3pt]
\midrule
\texttt{lem:grad-omega-split}
  & The orthogonal magnitude/direction split
  & \PROVED{} --- named by the audit as a survivor \\[3pt]
\texttt{eq:sigma-def}, its scale invariance
  & \(\sigma=A_{\mathrm{loc}}/M^{3/2}\)
  & \PROVED{} --- named by the audit as a survivor \\[3pt]
\midrule
\texttt{prop:coercive}
  & Coercivity of the second-derivative term
  & Rests on \texttt{lem:A-evol}; \OPEN \\[3pt]
\texttt{lem:A-evol}
  & \(-c\nu\int\abs{\nabla^2\omega}^2\varphi^2\) from diffusion
  & \OPEN: asserted without proof; \(\ehat\) is singular on
    \(\{\omega=0\}\) and the cross terms are sign-indefinite (D8) \\[3pt]
\texttt{thm:sigma-gronwall}
  & \(\int_0^TM\,dt<\infty\Rightarrow\sigma(T)<\infty\)
  & \CONDITIONAL{BKM}, hence useless as a route: the hypothesis is an
    equivalent of the conclusion (D5) \\[3pt]
\texttt{thm:blowup-alignment} (Theorem~E)
  & Alignment at blow-up; Type-I lower bound
  & \OPEN. Not among the five survivors the audit names; its proof invokes
    \texttt{thm:sigma-gronwall} (D5) and the Regime~I dichotomy, which needs
    \(\nu>C_1C/c\) (D6) \\[3pt]
\texttt{cor:log-BKM}
  & Log-corrected BKM criterion
  & \OPEN, same dependency \\[3pt]
The pressure-misalignment estimate
  & \(\abs{\nabla p_{\mathrm{glob}}/M}_{B_{\rstar}}\ge c_0\) for aligned tubes
  & \OPEN{} then, \OPEN{} now; never attempted after \Sn{31} \\
\bottomrule
\end{tabular}
\caption[Status of the Prize-claim early work]{Status of every object from
\Sn{29}--\Sn{30}. Two survive as proved; four are withdrawn; the remainder are
open, several of them resting on ingredients the audit found unproved.}
\label{tab:prize-status}
\end{table}

\begin{obsv}[The two survivors are the seed of everything after]
\label{obs:prize-survivors}
Of the thirteen objects in Table~\ref{tab:prize-status}, exactly two are
\PROVED: an algebraic identity and a scaling exponent. Both came from \Sn{30},
the session that produced no correct theorems. Both are used in every session
from \Sn{32} to \Sn{47}: the split identity is the pointwise device that makes
the \Sn{34} magnitude Caccioppoli work and the \Sn{36} Gram negativity
meaningful, and \(\sigma\)'s scale invariance is inherited by the localised
\(\sstar\) that the entire later programme is built on
(\texttt{prop:sigma-controls}: \(\sstar\le16\sigma\)).
Chapter~\ref{ch:inheritance} traces this in full.
\end{obsv}

\begin{obsv}[The failure mode, named]
\label{obs:prize-failure-mode}
The three fatal errors of \S20 --- D1, D2, D3 --- are all of one kind: an
inequality written down in the direction the argument needed, without the
computation that would have decided its direction being carried out. The
\emph{shape} of the argument was designed first and the estimates were fitted
to it. This is the failure mode the \Sn{41} merge gate
(Chapter~\ref{ch:house-rules}) exists to catch, and it is the same failure
mode named again in Wall~7 of the walls document, where a favourable term is
described as ``written down before the computation producing it was closed''.
It recurred at least three times after \Sn{31}, each time smaller, each time
caught earlier.
\end{obsv}
